\documentclass{amsart}

%%%%%%%packages
\usepackage{amssymb}
\usepackage{enumerate}
%\usepackage{showkeys}
%\usepackage{draftcopy}
\usepackage{xspace}
\usepackage{hyperref}

%%%%%%%%%Pagination
\hfuzz=15pt

%%%%%%%%%Theorems
\newtheorem{thm}{Teorema}[section]
\newtheorem{lem}[thm]{Lemma}
\newtheorem{cor}[thm]{Corollario} 
\newtheorem{sublem}[thm]{Sub-lemma}
\newtheorem{prop}[thm]{Proposizione}
\newtheorem{conj}{Congettura}\renewcommand{\theconj}{}
\newtheorem{defin}[thm]{Definitione}
\newtheorem{cond}{Conditione}
\newtheorem{exmp}[thm]{Esempio}
\newtheorem{es}[thm]{Esercizio}
\newtheorem{rem}[thm]{Commento}

%%%%%%%%%%mathcal
\newcommand\B{{\mathcal B}}
\newcommand\Co{{\mathcal C}}
\newcommand\D{{\mathcal D}}
\newcommand\F{{\mathcal F}}
\newcommand\Lp{{\mathcal L}}
\newcommand\M{{\mathcal M}}
\newcommand\Pa{{\mathcal P}}
\newcommand\T{{\mathcal T}}

%%%%%%%%%%%%mathbb
\newcommand\A{{\mathbb A}}
\newcommand\C{{\mathbb C}}
\newcommand\E{{\mathbb E}}
\newcommand\N{{\mathbb N}}
\newcommand\R{{\mathbb R}}
\newcommand\Q{{\mathbb Q}}
\newcommand\To{{\mathbb T}}
\newcommand\Z{{\mathbb Z}}

%%%%%%%%%%%greek
\newcommand\ve{\varepsilon}
\newcommand\vf{\varphi}



%%%%%   Figure-picture related definitions--beginning
\newenvironment{myfigure}{\begin{figure}[tb]\
\centering}{\end{figure}}

%\definecolor{lgray}{gray}{0.9}
%\definecolor{gray}{gray}{0.8}
%\definecolor{dgray}{gray}{0.7}
\setlength{\unitlength}{.8mm}
%%%%%   Figure related definitions--end


\pagestyle{myheadings}
\markboth{\centerline{\sc Carlangelo Liverani}\hskip-12pt}
{\centerline{\sc problemi concernenti i limiti }\hskip-12pt}

\begin{document}
\title{\bfseries Sul contare}
\author{Carlangelo Liverani}
\address{Dipartimento di Matematica\\
Universit\`{a} di Roma (Tor Vergata)\\
Via della Ricerca Scientifica, 00133 Roma, Italy\\
{\tt liverani@mat.uniroma2.it}}
\date{Rome, October 11, 2003}
\maketitle


\section{Un linguaggio}\label{sec:zero}
Una parte della diffolt\`a della matematica risiede nel fatto che essa
usa un linguaggio suo, in parte differente dal linguaggio naturale che
usiamo ogni giorno. In particolare, pu\`o capitare che parole del
linguaggio comune siano usate con una differente accezione in
matematica. Tuttavia ogni parola in matematica \`e precisamente
definita in termini di parole gi\`a note o dei termini {\em primitivi}
della teoria. Quando si legge un testo di matematica bisogna sempre
essere allerti alla possibilt\`a che una parola sia usata in senso
{\em tecnico} e in tal caso \`e buona norma domandarsi: che cosa
significa?


In questo corso si utilizzeranno
alcuni rudimenti del linguaggio logico  ed insiemistico. Poich\`e tale
elementi del linguaggio matematico sono di
uso abba\-stanza comune nelle scuole inferiori, essi saranno considerati
noti. Tuttavia, forse \`e meglio ricapitolare il minimo che verr\`a
usato nel corso. 

Per \emph{insieme} intender\`o una collezione di oggetti (o elementi),
poich\`e tutti gli 
insiemi trattati nel corso saranno alquanto concreti ed esenti da
problemi (tipicamente sottoinsiemi dei numeri reali) ignorer\`o
tranquillamente la distinzione tra insiemi e classi necessaria per non incorrere in
contraddizioni.\footnote{Essenzialmente le classi sono 
collezioni di oggetti e gli insiemi classi che appartengono ad altre
classi. Abbastanza sorprendentemente se si ammette che ogni classe pu\`o
appartenere ad un'altra classe (e quindi ogni classe \`e un insieme)
si incorre in una contraddizione.} 

Con $x\in A$ si intende
\emph{l'elemento $x$ appartiene all'insieme A}. Con $A\subset B$ si
intende \emph{l'insieme $A$ \`e contenuto nell'insieme $B$}, pi\'u
precisamente \emph{per ogni $x$ se $x\in A$ allora $x\in B$} o, pi\'u
sinteticamente, $\forall x(x \in A \rightarrow x\in B)$. 

\begin{exmp}
Se $A=\{1,2,3\}$ e $B=\{1,2,3,5\}$ allora $1\in A$ e $A\subset
B$. Attenzione per\`o: un insieme pu\`o anche essere un elemento di
un'altro insieme. Per esempio se $C=\{1,2,3, 5,\{1,2,3\}\}$ allora
$A\subset C$ ma anche $A\in C$ mentre $B\subset C$ ma $B\not\in C$,
cio\`e \emph{$B$ non \`e un elemento di $C$}.
\end{exmp}
 
Due insiemi $A,B$ si dicono uguali se $A\subset B$ e $B\subset
A$. Abbiamo poi l'unione di due insiemi $A\cup B=\{x: x\in A \text{
oppure } x\in B\}$, ovvero \emph{l'insieme di tutti gli elementi $x$
tali che $x$ appartiene ad $A$ oppure a $B$, o ad entrambi}, e
l'intersezione $A\cap B=\{x: x\in A \text{ 
e } x\in B\}$.\footnote{Si noti la possibilit\`a di definire un
insieme specificando una propriet\`a a cui i suoi elementi debbono
soddisfare. In realt\`a questo modo di definire un insieme pu\`o
generare paradossi se usato indiscriminatamente, per ovviare a questo
pericolo in generale occorrerebbe sviluppare una teoria assiomatica degli
insiemi (approssimativamente, il punto \`e che una propriet\`a
determina sempre una classe, che poi essa sia o meno un insieme \`e una cosa
che va dimostrata) che per\`o esula dagli scopi del presente corso,
anche perch\`e non veramente
necessaria per il tipo di propriet\`a a cui saremo interessati.} 
L'insieme vuoto $\emptyset$, cio\`e un buffo insieme che
non contiene alcun elemento.\footnote{L'esistenza di tale insieme,
come di altri che stiamo definendo, \`e in realt\`a un assioma nella
teoria assiomatica degli insiemi il che significa che piace tanto ai
matematici che non hanno alcuna intenzione di farne a meno.} Inoltre
abbiamo l'insieme potenza $\Pa(A)=\{x:x\subset A\}$, cio\`e
\emph{l'insieme di tutti i sottoinsiemi di $A$}.

Un altro modo di costruire nuovi insiemi partendo da insiemi dati \`e
il prodotto Cartesiano. Dati due insiemi $A$ e $B$ si dice loro
prodotto, e si scrive $A\times B$, l'insieme $\{(a,b):a\in A,\, b\in
B\}$. Per $(a,b)$ si intende la coppia ordinata $(a,b)$, cio\`e
$(a,b)$ \`e, in generale, diverso da $(b,a)$. In particolare, $A\times
B$ e $B\times A$ sono, in generali, diversi. Si noti che dati tre
insiemi, $A,B,C$,  $(A\times B)\times C$ e $A\times (B\times C)$ sono
naturalmente {\em isomorfi}\footnote{Questo significa che possono
essere posti in {\em corrispondenza biunivoca}, si veda pi\`u sotto per
una definizione precisa.} dove l'isomorfismo \`e stabilito dalla
corrispondenza $((a,b),c)\to (a,(b,c))$. \`E quindi naturale
sopprimere le parentesi. Dato un insieme $A$ si pu`o dunque farne ``il
quadrato'' cio\`e $A^2:=A\times A$ o, pi\'u in generale, l'ennesima
potenza, cio\`e $A^n:=A\times
\cdots \times A$, $n$ volte. Per ragioni che appariranno ovvie tra un
secondo si usa anche scrivere $A^{\{1,\dots,n\}}$.

Una funzione $F$ tra due insiemi $A$ e $B$ si scrive $F:A\to B$ ed \`e
semplicemente una regola tale che ad ogni $x\in A$ associa un unico elemento di
$B$ che viene appunto chiamato $F(x)\in B$. Una funzione si dice
\emph{suriettiva} se per ogni $y\in B$ esiste $x\in A$ tale che $y=F(x)$
(scritto brevemente $\forall y\in B \;\exists x\in A : y=F(x)$) cio\`e
ogni elemento di $B$ \`e l'immagine, attraverso $F$, di un qualche
elemento di $A$. Una funzione di dice \emph{iniettiva} se
$F(x)=F(y)\rightarrow x=y$ cio\`e un elemento di $B$ non pu\`o essere
l'immagine di due diversi elementi di $A$. L'insieme $A$ si dice
\emph{dominio} della funzione $F$ e $F(A)=\{F(x): x\in A\}$ si dice
\emph{codominio}. Finalmente, una funzione 
suriettiva ed iniettiva si dice \emph{biunivoca}. Se la funzione
$F:A\to B$ \`e biunivoca allora essa determina univocamente una
funzione $F^{-1}:B\to A$ tale che $\forall a\in A \;\; F^{-1}(F(a))=a$ e
$\forall b\in B \;\; F(F^{-1}(b))=b$. Tale funzione \`e detta
\emph{funzione inversa} ed \`e chiaramente una funzione biunivoca tra
$B$ ed $A$.

Considereremo inoltre dato (e noto) l'insieme (infinito) dei numeri naturali
$\N=\{0,1,2,\dots\}$.

\begin{lem}\label{lem:prod}
Per ogni insieme $A$ e $n\in\N$ gli insiemi $A^{\{1,\dots,n\}}$ e
$\{f:\{1,\dots,n\}\to A\}$ possono essere messi in corrispondenza
biunivoca. 
\end{lem}
\begin{proof}
Dato un elemento $\bar a=(a_1, \dots,a_n)\in A^n$ possiamo definire la
funzione $f_{\bar a}(i)=a_i$, $i\in\{1,\dots,n\}$. D'altro canto per ogni
funzione $f:\{1,\dots,n\}\to A$ possiamo definire la $n$-nupla ordinata
$a_f:=(f(1),\dots,f(n))$. Chiaramente, $a_{f_{\bar a}}=\bar
a$. Possiamo dunque definire la corrispondenza 
biunivoca $F:A^n\to \{f:\{1,\dots,n\}\to A\}$, $F(\bar a)=f_{\bar a}$.
\end{proof}

Tramite la funzione $F$ costruita nella dimostrazione del Lemma
\ref{lem:prod} possiamo dunque identificare i due
insiemi. Abbiamo cos\'\i\  una definizione alternativa di
$A^{\{1,\dots,n\}}$. Il vantaggio di tale definizione \`e che ora
$A^B$ \`e ben definito per ogni insieme $A$ e $B$, cio\`e
$A^B=\{f:B\to A\}$.

 
Gli insiemi possono essere finiti o infiniti, tuttavia per dare un
significato preciso a questa affermazione occorre dire che si intende
per \emph{finito} o \emph{infinito}. A questo scopo \`e utile chiarire
che significa dire che due insiemi hanno \emph{lo stesso numero di elementi}.

\begin{defin}\label{def:cont}
Diremo che due insiemi $A$ e $B$ hanno lo stesso numero di elementi o
che hanno \emph{la stessa cardinalit\`a} se esiste una funzione biunivoca
$F:A\to B$.\footnote{Infatti, se ne esiste una, di solito ne esisteranno
molte altre.}
\end{defin}

\section{uno due tre ...}\label{sec:uno}
Cerchiamo di familiarizzarci con le conseguenze della definizione
\ref{def:cont}.

\begin{es}\label{es:uno}
Si mostri che se $A$ e $B$ hanno la stesa cardinalit\`a e se $B$ e $C$
hanno la stessa cardinalit\`a allora $A$ e $C$ hanno la stessa cardinalit\`a.
\end{es}

\begin{defin}\label{def:enum}
Diremo che un insieme ha $n\in\N$ elementi se ha la medesima cardinalit\`a
dell'insieme $\{1,2,\dots,n\}\subset\N$.\footnote{Ovviamente, diremo
che l'insieme vuoto ha zero elementi.}
\end{defin}
Tale definizione \`e sensata, ma per
accertarcene occorre fare il seguente esercizio che richiede un attimo
di meditazione.
\begin{es}\label{es:consi}
Si mostri che un insieme non pu\`o avere sia $n$ che $m$ elementi, per
$n\neq m$. (Suggerimento: per assurdo. Se non fosse vero esisterebbero
$n,m\in\N$, $n>m$, tali che gli insiemi $\{1,2,\dots,n\}$ e
$\{1,2,\dots,m\}$ sono in corrispondenza biunivoca. Da questo si deduca
che lo stesso deve valere per $n-1$ e $m-1$. Per fare questo si mostri
che se un insieme ha $n$ elementi e se ne toglie uno allora l'insieme
rimanente ha $n-1$ elementi.\footnote{Pu\'o sembrare pazzesco che si
debba dimostrare una cosa del genere, ma bisogna ricordare che, in
questo contesto,
\emph{avere $n$ elementi} non ha necessariamente nulla a che fare col
solito significato intuitivo ma ha il preciso significato specificato
nella definizione \ref{def:enum}. Che questo significato sia
consistente col significato intuitivo che abbiamo di questo concetto
\`e esattamente quello che stiamo verificando in questo momento.} 
Iterando l'argomento ci
si riduce a $n-m+1>1$ e $1$, cosa assurda.) 
\end{es}
Fino ad ora tutto sembra alquanto banale, anche se la spiegazione di
che si intende per contare pu\`o sembrare piuttosto strana, tuttavia
sembra coincidere con la nostra intuizione. Va comunque notato che si possono avere
problemi un poco pi\'u interessanti di quello dell'esercizio \ref{es:uno}.
\begin{lem}\label{lem:potfun}
Per ogni insieme $A$, gli insiemi $\Pa(A)$ e $\{0,1\}^A$ hanno lo
stesso numero di elementi.
\end{lem}
\begin{proof}
Definiamo la funzione $F:\Pa(A)\to \{0,1\}^A$ nel modo seguente: per
ogni $B\subset A$ sia $F(B)$ la funzione che vale $1$ in $B$ e zero
fuori, cio\`e, per ogni $a\in A$,
\[
F(B)(a)=\left\{\begin{aligned}
	       &1\quad\text{se }a\in B\\
	       &0\quad\text{se }a\not\in B
               \end{aligned}
\right.
\]
A volte $F(B)$ \`e chiamata la \emph{funzione caratteristica} di
$B$. Chiaramente $F$ \`e biunivoca e questo conclude la dimostrazione.
\end{proof}

Tra le altre cose il lemma precedente chiarisce perch\`e $\Pa(A)$ \`e
chiamato l'insieme \emph{potenza}.

Tuttavia, le vere sorprese giungono
quando si cerca di applicare tali concetti agli insiemi infiniti (di
cui al momento abbiamo una sola istanza: $\N$).
\begin{exmp}\label{ex:2}
Si consideri l'insieme $\N_*=\{1,2,\dots\}\subset\N$. Tale insieme ha
un elemento in meno dei naturali (lo zero) e quindi ci si potrebbe
aspettare che la sua cardinalit\`a sia inferiore. Invece se si
considera $F:\N_*\to\N$ definita da $F(n)=n-1$ si ha chiaramente un
funzione biunivoca tra $\N_*$ e $\N$. Dunque, per definizione, $\N_*$ e
$\N$ hanno lo stesso numero di elementi! 
\end{exmp}

\begin{exmp}\label{ex:3}
Si consideri l'insieme $\N_p=\{0,2,4,\dots, 2n,\dots\}$ dei numeri
pari. Di nuovo si potrebbe pensare che contenga la met\`a degli
elementi di $\N$, cio\`e che la sua cardinalit\`a sia
inferiore. Eppure, applicando la definizione e considerando la funzione
$F(n)=2n$ si ottiene che $\N_p$ e $\N$ hanno la stesa cardinalit\`a.
\end{exmp}
Queste stranezze hanno a lungo lasciato perplessi i matematici, fino a
quando non ci si \`e rassegnati alla constatazione che l'infinito \`e
un concetto controintuitivo e si sono usate le sue peculiarit\`a
proprio per definirlo.

\begin{defin}\label{def:infy}
Un insieme $A$ si dice \emph{infinito}, se esiste un suo sottoinsieme
proprio (cio\`e un $B\subset A$ ma $B\neq A$) con la stessa
cardinalit\`a di $A$.
\end{defin}
Vediamo che una tale definizione non fa danni.
\begin{es}\label{es:ff}
Se $A$ \`e un insieme finito (cio\`e non \`e infinito) allora ogni suo
sottoinsieme \`e finito. (Suggerimento: per assurdo. Si
supponga che esistano $C\subset B\subset A$ tali che $C$ pu\`o essere messo in
corrispondenza biunivoca con $B$ e si costruisca un sottoinsieme $C_1$
che pu`o essere messo in corrispondenza con $A$.)
\end{es}

\begin{es}\label{es:finito}
Per ogni $n\in\N$ l'insieme $A_n:=\{0,1,\dots,n\}$ \`e
finito. (Suggerimento: per assurdo si supponga che esista $n\in \N$
tale che $A_n$ \`e infinito. Sia $B\subset A_n$ in corrispondenza
biunivoca con $A_n$ tramite la funzione $F$. Si tolga da $A$ il pi\'u
grande elemento che non \`e contenuto in $B$, diciamo $k$, e si tolga
da $B$ l'elemento $F^{-1}(k)$. Si mostri che in tal modo segue che anche
$A_{n-1}$ \`e infinito. Continuando in tale modo si ottiene che $A_0$
\`e in corrispondenza biunivoca con un suo sottoinsieme proprio,
insensato visto che l'unico sottoinsieme proprio \`e $\emptyset$.)
\end{es}
Dunque con la definizione \ref{def:infy} gli insiemi finiti sono
quelli che ci aspettiamo mentre quelli infiniti godono delle strane
patologie di cui agli esempi \ref{ex:2}, \ref{ex:3} 
In realt\`a, patologie anche peggiori sono possibili.
\begin{es}\label{es:raz} Si mostri che gli insiemi $\N$ ed $\N\times
\N_*$ hanno lo stesso numero di elementi. (Suggerimento: si provi con la funzione
$F:\N\times \N_*\to\N$ definita da
$F((p,q))=p+1+\sum_{i=0}^{p+q}(i+1)$).\footnote{La funzione \`e
ottenuta considerando gli insiemi $G_n=\{(p,q)\in\N\times \N_* :
p+q=n\}$, chiaramente $\cup_{n\in\N}G_n=\N\times \N_*$. A questo punto si contano
prima gli elementi in $G_0$, poi in $G_1$ e cosi via, in ogni $G_i$, invece,
gli elementi sono contati secondo l'ordine di $p$ crescente.}
\end{es}
Si noti che l'insieme dei numeri razionali pu\`o essere considerato
come contenuto in $\N\times \N_*$, infatti ad ogni $(p,q)\in\N\times
\N_*$ si pu\`o associare il numero razionale $\frac pq$. Ovviamente
tale funzione non \`e iniettiva ma questo significa solo che il numero
degli elementi di $\N\times \N_*$ \`e maggiore od uguale a quello di
$\Q$, ma poich\`e $\N\subset \Q$ \`e chiaro che $\N$ e $\Q$ hanno lo
stesso numero di elementi.\footnote{Ovviamente, dire che un insieme
$A$ ha una cardinalit\`a maggiore o uguale di un insieme $B$ significa
semplicemente che esiste $C\subset A$ con la stessa cardinalit\`a di $B$.}

A questo punto si potrebbe pensare che tutti gli insiemi infiniti hanno
la stessa cardinalit\`a e dunque, in un qualche senso, esiste un solo
tipo di infinito. Anche questo risulta essere falso.
\begin{lem}\label{lem:poti}
Per ogni insieme $A$ la cardinalit\`a di $\Pa(A)$ \`e sempre maggiore
della cardinalit\`a di $A$.
\end{lem}
\begin{proof}
Ragionando per assurdo, supponiamo che il lemma sia falso e che quindi
esista una funzione 
biunivoca $F:A\to\Pa(A)$. In tal caso definiamo l'insieme $B=\{x\in A:
x\not\in F(x)\}$. Chiaramente $B\subset A$ e dunque
$B\in\Pa(A)$. Esister\`a perci\'o un unico elemento $z\in A$ tale che
$F(z)=B$. A questo punto possiamo ottenere una contraddizione
chiedendoci se $z\in B$ oppure no. Infatti se $z\in B$ allora, per
definizione di $B$, $z\not\in F(z)=B$. Dunque deve essere $z\not\in
B$. Ma se $z\not\in B$ allora $z\not\in F(z)$ e, per la definizione di
$B$, $z\in B$, assurdo.
\end{proof}
Ci\`o significa che il numero degli elementi di $\Pa(\N)$ \`e
strettamente maggiore di quello di $\N$. Si noti che $\{0,1\}^\N$
pu\`o essere interpretato come la collezione di stringhe infinite
$(a_0,a_1,\dots, a_n,\dots)$, dove $a_i\in\{0,1\}$ per ogni $i\in\N$. 
Ognuna di tali stringhe pu`o essere pensata come l'espansione binaria
di un numero reale nell'intervallo $[0,1]$, tuttavia diverse stringhe
possono corrispondere allo stesso numero (ad esempio $(1,0,0,\dots)$ e
$(0,1,1,1,\dots)$ corrispondono entrambi al numero uno, come vedremo
nel proseguo del corso). Si pu\`o comunque dimostrare che solo un
numero numerabile di stringhe sono problematiche e dunque la
cardinalit\`a di $[0,1]$  coincide con
quella di $\Pa(\N)$.

Per concludere vorrei mettere in guardia contro la possibile
impressione che contare insiemi finiti sia una banalit\`a. In realt\`a
contare pu\`o essere un problema assai difficile,\footnote{\`E il
soggetto di una intera branca della matematica: la combinatorica.} di
solito occorre una qualche idea molto astuta su come organizzare il conteggio per
venirne a capo. Allo scopo di capire meglio il problema facciamoci una
semplice domanda: per gli insiemi finiti quanto \`e grande l'insieme potenza?
\begin{lem}\label{lem:potf}
Per ogni insieme $A$ di cardinalit\`a $n$, l'insieme
$\Pa(A)$ ha cardinalit\`a $2^n$.
\end{lem}
\begin{proof}
Il lemma \`e banale per $A=\emptyset$. Per le altre posibilt\`a useremo
il risultato del Lemma \ref{lem:potfun} e cercheremo di contare gli
elementi dell'insieme $\{0,1\}^A$. Organizziamo il conteggio attraverso
la seguente identit\`a algebrica: 
per ogni due numeri $\alpha_0, \alpha_1$ si ha\footnote{Dimostrarla
equivale a capire cosa si \`e scritto ed \`e un esercizio molto
utile. La si dimostri prima per $A=\{1,2,\dots,n\}$ e poi per un
generico insieme $A$ con $n$ elementi. Per la cronaca $\sum\limits_{i\in
\{1,\dots,n\}}\beta_i:=\beta_1+\beta_2 +\cdots+\beta_n$ e $\prod\limits_{i\in
\{1,\dots,n\}}\beta_i:=\beta_1\beta_2 \cdots\beta_n$.}
\begin{equation}\label{eq:compl}
(\alpha_0+\alpha_1)^n=\sum_{f\in\{0,1\}^A}\;\prod_{i\in A}
\alpha_{f(i)}. 
\end{equation}
Utilizzando tale formula con $\alpha_0=\alpha_1=1$ si ottiene
\[
2^n=(1+1)^n=\sum_{f\in\{0,1\}^A}\;\prod_{i\in A} 1
=\sum_{f\in\{0,1\}^A} 1. 
\]
Poich\`e l'ultima espressione \`e semplicemente la somma di uno per
ogni elemento di $\{0,1\}^A$ \`e chiaro che la somma da esattamente la
cardinalit\`a di $\{0,1\}^A$ e dunque di $\Pa(A)$.
\end{proof}

A qualcuno la formula \eqref{eq:compl} pu\`o sembrare assai complessa,
ci si pu\`o chiedere se \`e possibile scriverla in maniera pi\'u
semplice. La risposta \`e positiva ma richiede di risolvere un altro
problema di conteggio.

Sia $B_k=\{f\in\{0,1\}^A :\sum\limits_{i\in
A}f(i)=k\}$,\footnote{Cio\`e $B_k$ \`e l'insieme di tutti i
sottoinsiemi di $A$ con cardinalit\`a $k$.} chiaramente
$\{0,1\}^A=\cup_{k=0}^n B_k$, dunque possiamo scrivere\footnote{Dato
un insieme finito $A$, per $\# A$ si intende il numero di elementi di $A$.}
\[
(\alpha_0+\alpha_1)^n=\sum_{k=0}^n\sum_{f\in B_k}\;\prod_{i\in A}
\alpha_{f(i)}=\sum_{k=0}^n\sum_{f\in B_k} \alpha_1^n\alpha_0^{n-k}
=\sum_{k=0}^n\# B_k \alpha_1^n\alpha_0^{n-k}.
\]
\begin{es}[Binomio di Newton]
Per ogni due numeri $a,b$ vale\footnote{Per $\binom nk$ si intende
$\frac{n!}{k!(n-k)!}$ mentre $k!=1\cdot 2\cdot 3\cdots n=\prod_{i=1}^n
i$.}
\[
(a+b)^n=\sum_{k=0}^n\binom nk a^nb^{n-k}.
\]
(Suggerimento: per contare gli elementi di $B_k$  si devono contare
tutti i possibili modi in cui si possono scegliere $k$ elementi tra
$n$ (sono gli elementi su cui $f$ varr\'a uno). Per farlo, prima si sceglie
chi \`e il primo elemento ($n$ possibilit\`a) poi chi \`e il secondo
($n-1$ possibilit\`a, visto che un elemento \`e gi`a stato scelto e
non si pu\`o scegliere nuovamente), e cos\'\i\  via. Questo da
$n(n-1)\cdots (n-k+1)=\frac {n!}{(n-k)!}$ possibilt\`a. Per\`o,
cos\`\i\  facendo, abbiamo scelto lo stesso insieme un mucchio di
volte, infatti se lo stesso insieme viene scelto in un ordine diverso
nulla cambia, mentre noi lo abbiamo contato di nuovo. Quante volte
abbiamo contato lo stesso insieme? Ovviamente $k!$.)
\end{es}

\end{document}

